	\documentclass[12pt]{exam}
	\usepackage[utf8]{inputenc}
	\usepackage[portuguese,brazil]{babel}
	\usepackage[margin=.65in]{geometry}
	\usepackage{amsmath,amssymb}
	\usepackage{multicol}
	\usepackage{multirow}
	\usepackage{float}
	\usepackage{enumerate}
	\usepackage{enumitem}
	\usepackage{array}
	\usepackage{graphicx}
	\usepackage{listings}
	\usepackage{circuitikz}    
	\usepackage{xcolor}
	\usepackage{tikz}
	\usetikzlibrary{shapes.multipart, calc}
	\usepackage{rotating}
	\usepackage{caption}
	%\usepackage[version=3]{mhchem}

	\pointpoints{ponto}{pontos}
	\hpword{Pontos:}
	\vpword{Pontos}
	\htword{Total}
	\vtword{Total:}
	\vsword{Resultado}

	\definecolor{azul}{RGB}{0,32,96}
	\definecolor{verde}{RGB}{0,132,124}
	\definecolor{cinza}{RGB}{0,84,132}
	%\definecolor{cinza2}{RGB}{218,220,216}
	\definecolor{cinza2}{RGB}{132,134,136}
	\definecolor{mygreen}{rgb}{0,0.6,0}
	\definecolor{mygray}{rgb}{0.5,0.5,0.5}
	\definecolor{mymauve}{rgb}{0.58,0,0.82}
	\newcommand*{\vttfamily}{%
	\fontencoding{T1}\fontfamily{lmvtt}\selectfont
	}
	\lstdefinestyle{Cstyle}{%
	language=C,numbers=left, keywordstyle=\color{mymauve}\vttfamily,commentstyle=\color{mygreen}\vttfamily,stringstyle=\color{blue},basicstyle=\scriptsize\vttfamily,numberstyle=\tiny\vttfamily,breaklines=true,stepnumber=1,tabsize=2,xleftmargin=\leftmarginii,escapechar=!,captionpos=b,morekeywords={include, printf,NULL},columns=flexible,escapeinside={(^}{^)},
	frame=l,
	framesep=2.5mm,
	framexleftmargin=2.5mm,
	fillcolor=\color{cinza2!50},
	rulecolor=\color{ballblue},
	rulecolor=\color{blue}
	}

	\lstnewenvironment{C}[1][]
	{
	\lstset{style=Cstyle,#1}
	}
	{}


	\newcommand{\class}{Segurança da Informação}
	\newcommand{\profname}{Gilmar Gomes do Nascimento}
	\newcommand{\examnum}{Exercício 0} \newcommand{\examdate}{$\today$}
	\newcommand*{\oldneg}{\mathord{\nao}}
	\pagestyle{headandfoot}
	\runningheadrule
	\firstpageheader{}{}{}
	\runningheader{\class}{\examnum}{\examdate}
	\runningfootrule
	\firstpagefooter{}{}{}
	\runningfooter{}{Página \thepage\ de \numpages}{}
	%\pagestyle{foot}
	%\lfoot{}
	%\cfoot[]{Página \thepage\ de \numpages}
	%\rfoot{}
	%
	%\runningfootrule
	%\boxedpoints %\pointsinmargin
	%\printanswers
	%\SolutionEmphasis{\color{red}}
	%\definecolor{SolutionBoxColor}{gray}{0.8}
	%\renewcommand{\solutiontitle}{\noindent\textbf{Solução:}\enspace}
	%\renewcommand{\solutiontitle}{\noindent\textbf{Solution:}\par\noindent}

	\begin{document}

	\noindent
	\begin{tabular*}{\textwidth}{l @{\extracolsep{\fill}} r @{\extracolsep{6pt}} l}
	\textbf{Disciplina:} \class &  \multicolumn{2}{c}{\multirow{4}{*}{
	\includegraphics[scale=0.35]{logo.png}}}\\
	\textbf{Prof.:} ~\profname \vspace{0.2in} && \\
	\vspace{0.2in} \textbf{Nome:} \makebox[3in]{\hrulefill} \vspace{0.2in} &&\\
	%\textbf{Matrícula:} \makebox[1.5in]{\hrulefill} \textbf{Nota:} \makebox[1in]{\hrulefill} &&\\ 
	%\multicolumn{3}{c}{\textbf{{\examnum \hspace{3px} \examdate}}}\\
	\multicolumn{3}{c}{}\\
	\end{tabular*}\\
	\rule[2ex]{\textwidth}{1.6pt}
	\parskip=0.1in

	%Este exame contêm \numpages\ páginas (incluindo esta capa) e \numquestions\ questões.\\
	%Total de \numpoints~ pontos.

	%\begin{center}
	%Grade Table (for teacher use only)\\
	%\addpoints
	%\gradetable[v][questions]
	%\end{center}

	\noindent

	\begin{questions}
		\question Descreva o conceito de SSRF (Server-Side Request Forgery).
			\begin{solutionbox}{2in} \end{solutionbox}
		\question Diferenças entre SSRF in-band e cego
			\begin{solutionbox}{2in} \end{solutionbox}

			\begin{solutionbox}{2.44in} \end{solutionbox}
		\question Escreva V ou F para cada assertiva sobre vulnerabilidades em SSRF:

			(\hspace(3px)) O navegador do usuário é enganado  \\
			(\hspace(3px)) O servidor é quem faz a requisição \\
			(\hspace(3px)) O alvo típico é a rede interna, localhost e APIs\\
			(\hspace(3px)) O SSRF é menos perigoso, pois as requisições se originam de um servidor confiável \\
			(\hspace(3px)) O alvo canônico para ataques de SSRF em ambientes de nuvem é: \textt{/latest/meta-data/iam/security-credentials/}\\
			(\hspace(3px)) Se acessível via SSRF, pode revelar credenciais temporárias (IAM roles) que permitem
			controle total sobre a infraestrutura \\
			(\hspace(3px)) APIs que consomem URLs externas (geradores de PDF, webhooks, importadores de imagem)
			podem se tornar vetores de ataque para SSRF. \\
			(\hspace(3px)) Qualquer funcionalidade que, por design, precise acessar uma URL fornecida pelo usuário
			é um vetor de SSRF em potencial.

		\question Discorra sobre \textit{webhooks} e correlacione com SSRF. 
			\begin{solutionbox}{2in} \end{solutionbox}

		
		\question O que são IAM roles? Pode ser utilizado para roubar dados com SSRF?
			\begin{solutionbox}{2.5in} \end{solutionbox}

		\question Correlacione as colunas
			\begin{multicols}{2}
				\begin{column}
					\begin{enumerate}
						\item Transferência de arquivos e bypass de firewall
						\item Acesso a dicionários e info gathering
						\item Comunicação com serviços internos
						\item Leitura de arquivos locais
					\end{enumerate}
				\end{column}
				\begin{column}
					\begin{itemize}
						\item (\hspace(3px)) dict
						\item (\hspace(3px)) ftp
						\item (\hspace(3px)) gopher
						\item (\hspace(3px)) file
					\end{itemize}
				\end{column}
			\end{multicols*}
		

		

	\end{questions}
	\end{document}